\documentclass[11pt]{beamer}
\usepackage[utf8]{inputenc}
\usepackage{amsmath, amssymb, graphicx, listings, xcolor}
\usepackage{bm}

% Presentation Theme
\usetheme{Madrid}
\usecolortheme{default}

% Code Block Styling
\lstset{
    language=Python,
    basicstyle=\ttfamily\scriptsize,
    keywordstyle=\color{blue}\bfseries,
    commentstyle=\color{green!50!black},
    frame=single,
    backgroundcolor=\color{gray!10}
}

\title[Earth to Jupiter Mission Design]{Astrodynamics 2026-1}
\subtitle{Capstone Mission: Earth to Jupiter \\ Analytical, Numerical, \& GMAT Solutions}
\author{Prof. Juan}
\institute{Universidad de Antioquia - Aerospace Engineering}
\date{Spring 2026}

\begin{document}

% --- Title Slide ---
\begin{frame}
    \titlepage
\end{frame}

% ==========================================
% SECTION 1: ANALYTICAL SOLUTION
% ==========================================
\section{Analytical Solution (Patched Conics)}

% --- Slide 1 ---
\begin{frame}{Analytical Assumption: Ideal Hohmann}
    We begin with the Patched Conic approximation, assuming circular, coplanar planetary orbits. Hohmann transfers are the easiest to analyze and the most energy efficient.
    
    \vspace{0.2cm}
    \textbf{Given Planetary Constants:}
    \begin{itemize}
        \item $\mu_{sun} = 1.327 \times 10^{11} \text{ km}^3/\text{s}^2$
        \item $R_{\oplus} = 149.6 \times 10^6 \text{ km}$ (Earth)
        \item $R_{\jupiter} \approx 5.203 \text{ AU} = 778.4 \times 10^6 \text{ km}$ (Jupiter)
    \end{itemize}
    
    \textbf{1. Heliocentric Transfer Orbit ($a_{tx}$):}
    \begin{equation}
        a_{tx} = \frac{R_{\oplus} + R_{\jupiter}}{2} = 464.0 \times 10^6 \text{ km}
    \end{equation}
    
    \textbf{2. Transfer Velocity at Earth ($V_{D}^{(v)}$):}
    \begin{equation}
        V_{D}^{(v)} = \sqrt{2\mu_{sun}}\sqrt{\frac{R_{\jupiter}}{R_{\oplus}(R_{\oplus}+R_{\jupiter})}} = \mathbf{38.56 \text{ km/s}}
    \end{equation}
\end{frame}

% --- Slide 2 ---
\begin{frame}{Analytical: Departure Energy ($C_3$)}
    To find the energy required from our launch vehicle, we subtract Earth's velocity to find the Hyperbolic Excess Speed ($v_\infty$).
    
    \vspace{0.2cm}
    \textbf{3. Earth's Heliocentric Velocity ($V_{\oplus}$):}
    \begin{equation}
        V_{\oplus} = \sqrt{\frac{\mu_{sun}}{R_{\oplus}}} = \mathbf{29.78 \text{ km/s}}
    \end{equation}
    
    \textbf{4. Hyperbolic Excess Speed ($v_\infty$):}
    \begin{equation}
        v_\infty = V_{D}^{(v)} - V_{\oplus} = 38.56 - 29.78 = \mathbf{8.78 \text{ km/s}}
    \end{equation}
    
    \textbf{5. Characteristic Energy ($C_3$):}
    \begin{equation}
        C_3 = v_\infty^2 = \mathbf{77.09 \text{ km}^2/\text{s}^2}
    \end{equation}
\end{frame}

% --- Slide 3 ---
\begin{frame}{Analytical: LEO Parking Orbit $\Delta v$}
    Assume our spacecraft is parked in a 300 km altitude Low Earth Orbit ($r_{p} = 6678 \text{ km}$). $\mu_{\oplus} = 398,600 \text{ km}^3/\text{s}^2$.
    
    \vspace{0.2cm}
    \textbf{6. Velocity in Parking Orbit ($v_c$):}
    \begin{equation}
        v_c = \sqrt{\frac{\mu_{\oplus}}{r_{p}}} = \mathbf{7.73 \text{ km/s}}
    \end{equation}
    
    \textbf{7. Required Velocity at Periapsis ($v_p$):}
    \begin{equation}
        v_p = \sqrt{v_\infty^2 + \frac{2\mu_{\oplus}}{r_{p}}} = \sqrt{77.09 + 119.38} = \mathbf{14.01 \text{ km/s}}
    \end{equation}
    
    \textbf{8. Trans-Jovian Injection Engine Burn ($\Delta v$):}
    \begin{equation}
        \Delta v = v_p - v_c = 14.01 - 7.73 = \mathbf{6.28 \text{ km/s}}
    \end{equation}
\end{frame}

% ==========================================
% SECTION 2: NUMERICAL SOLUTION
% ==========================================
\section{Numerical Solution (SPICE \& Lambert)}

% --- Slide 4 ---
\begin{frame}{The Reality of the Jovian System}
    Our analytical math yielded $\Delta v = 6.28 \text{ km/s}$. However, planets' orbits are not perfectly coplanar, and the transfer ellipse is often tangent to neither orbit. Jupiter has an orbital inclination of $1.3^\circ$ and eccentricity of $0.048$.
    
    \vspace{0.2cm}
    \textbf{Targeting Jupiter numerically:}
    \begin{itemize}
        \item Non-Hohmann transfers require the solution of Lambert's problem.
        \item We obtain the state vector of Earth ($R_1$) at departure and Jupiter ($R_2$) at arrival.
        \item Lambert's problem yields the required departure velocity $V_{D}^{(v)}$ relative to the sun.
    \end{itemize}
    
    \vspace{0.2cm}
    To find a \textit{real} mission window, we use a Python script to sweep through thousands of Lambert solutions to generate a Porkchop Plot.
\end{frame}

% ==========================================
% SECTION 3: GENERATING THE PORKCHOP PLOT
% ==========================================
\section{Generating the Porkchop Plot}

% --- Slide 5 ---
\begin{frame}[fragile]{Python/SPICE: Setting up the Jupiter Grid}
    First, we adjust our Colab script to target Jupiter (NAIF ID: 599) and define our time grid based on the Earth-Jupiter synodic period (approx. 398 days).
    
\begin{lstlisting}
# Constants & Target
EARTH = '399'
TARGET_ID = '599' # Jupiter Barycenter

# Define the search windows for a specific opportunity
dep_start = spy.str2et('2028-10-01')
dep_end = spy.str2et('2028-12-01')

arr_start = spy.str2et('2031-01-01')
arr_end = spy.str2et('2032-06-01')

# Initialize 2D grid to store C3 values
c3_grid = np.full((len(arrival_dates), len(departure_dates)), np.nan)
\end{lstlisting}
\end{frame}

% --- Slide 6 ---
\begin{frame}[fragile]{Python/SPICE: Populating the $C_3$ Matrix}
    We iterate over every departure and arrival date, solve Lambert's problem, and populate our 2D grid with the required Characteristic Energy ($C_3$).
    
\begin{lstlisting}
for i, t_dep in enumerate(departure_dates):
    r1, v_earth = get_state(EARTH, t_dep)
    
    for j, t_arr in enumerate(arrival_dates):
        r2, _ = get_state(TARGET_ID, t_arr)
        tof = t_arr - t_dep
        
        # Solve Lambert's Problem
        solutions = list(izzo.lambert(MU_SUN, r1, r2, tof))
        v_tx1, v_tx2 = solutions[0].value 
        
        # Calculate C3 and save to matrix
        v_inf_vec = v_tx1 - v_earth
        c3_grid[j, i] = np.dot(v_inf_vec, v_inf_vec)
\end{lstlisting}
\end{frame}

% --- Slide 7 ---
\begin{frame}{Analyzing the Earth-Jupiter Porkchop Plot}
    When we plot \texttt{c3\_grid} using \texttt{matplotlib.contourf}, the topography of the launch window emerges.
    
    \vspace{0.3cm}
    \textbf{What to look for on the Jupiter plot:}
    \begin{itemize}
        \item \textbf{High Energy Ridge ($\Delta \nu = \pi$):} A massive energy spike dividing the plot. Because Jupiter's orbit is inclined, a perfect 180-degree transfer requires an impossible plane-change maneuver.
        \item \textbf{Type I Trajectory Lobe:} Shorter time of flight ($< 180^\circ$ around the sun). Often requires a slightly higher $C_3$.
        \item \textbf{Type II Trajectory Lobe:} Longer time of flight ($> 180^\circ$).
        \item \textbf{The Optimal Valley:} Because of 3D realities, the absolute minimum $C_3$ valley usually sits between \textbf{$80 \text{ to } 85 \text{ km}^2/\text{s}^2$}. This translates to an actual $\Delta v$ closer to \textbf{$6.5 \text{ km/s}$} (higher than our 6.28 km/s ideal analytical guess).
    \end{itemize}
\end{frame}

% ==========================================
% SECTION 4: GMAT VALIDATION
% ==========================================
\section{GMAT Validation}

% --- Slide 8 ---
\begin{frame}[fragile]{NASA GMAT: Mission Setup}
    Once our script finds the minimum $C_3$ valley, we extract the optimal $\mathbf{v}_{\infty}$ vector and input it into GMAT's Differential Corrector.
    
    \vspace{0.3cm}
    \textbf{The Targeting Script:}
\begin{lstlisting}
% 1. The Control Variables (Initial Guesses)
% Use our numerical Delta V (~6.5 km/s) and Beta angle
Vary DC1(TJI_Burn.Element1 = 6.5);
Vary DC1(DefaultSC.TA = 120.0);
Vary DC1(DefaultSC.RAAN = 0.0);

% 2. Propagate Outward
Maneuver TJI_Burn(DefaultSC);
Propagate DefaultProp(DefaultSC) {DefaultSC.Earth.RMAG = 1e6};

% 3. The Targets (From Python SPICE Output)
Achieve DC1(DefaultSC.Earth.C3Energy = 82.5);
Achieve DC1(DefaultSC.Earth.RLA = -15.4);
Achieve DC1(DefaultSC.Earth.DLA = 5.2);
\end{lstlisting}
\end{frame}

% --- Slide 9 ---
\begin{frame}{GMAT: Arrival at Jupiter}
    If the Departure sequence converges, the spacecraft is on a ballistic path to Jupiter. We use a second \texttt{Propagate} command to coast across the solar system.
    
    \vspace{0.3cm}
    \textbf{The B-Plane Targeting Sequence:}
    To hit a specific Jovian orbit (like the Galileo mission, which arrived at Jupiter in December 1995), we must use a second Differential Corrector midway through the cruise phase to target Jupiter's \textbf{B-Plane}.
    
    \begin{itemize}
        \item \textbf{BdotT:} Targets the radial distance/altitude.
        \item \textbf{BdotR:} Targets the orbital inclination upon arrival.
    \end{itemize}
    
    \vspace{0.3cm}
    \textit{This multi-step approach—Analytical Estimate $\rightarrow$ Numerical Lambert Optimization $\rightarrow$ High-Fidelity GMAT targeting—is exactly how modern interplanetary missions are flown!}
\end{frame}

\end{document}